\documentclass[12pt, a4paper]{article}
\usepackage[utf8]{inputenc}
\usepackage[brazilian]{babel}
\usepackage[T1]{fontenc}
\usepackage{graphicx}
\usepackage{geometry}
\usepackage{color,soul}
\DeclareRobustCommand{\hlcyan}[1]{{\sethlcolor{cyan}\hl{#1}}}
\usepackage[document]{ragged2e}

\newcommand{\+}[1]{\ensuremath{\mathbf{#1}}} % comando para escrever matrizes em negrito \+A= A

\geometry{top=2.5cm, bottom=2.5cm, right=2.5cm, left=2.5cm}

\begin{document}
\noindent Universidade Federal Rural de Pernambuco \hfill Teoria do Desenvolvimento Rural\\
\noindent PADR \hfill 2019\\
\noindent Prof. Diego Firmino Costa da Silva \\
\noindent Aluno: Nome completo do aluno \\ % Inserir o nome completo do aluno
\center{Resenha do artigo:}
\center{\large{Título do artigo}} % Inserir o título do artigo
\center{\textit{Autores do artigo}} % Inserir o nome dos autores dos artigos
\justify

%%%%%%% Conteúdo da resenha %%%%%%%%%%%

Como uma síntese, a resenha deve ir direto ao ponto, mesclando momentos de pura descrição com momentos de crítica direta. O resenhista que conseguir equilibrar perfeitamente esses dois pontos terá escrito a resenha ideal.

A resenha deve possuir as mesmas qualidades de estilo imprescindíveis a todo texto escrito, como: simplicidade, clareza, concisão, propriedade vocabular, precisão vocabular, objetividade e impessoalidade. Além disso, a resenha deve apresentar imparcialidade, atitude científica e privilegiar o essencial.

Em relação à Seja na defesa ou no ataque, o resenhista deve julgar as ideias da obra sem paixão, devendo posicionar-se criticamente como um juiz e apresentar tanto os aspectos positivos quanto os negativos da obra. A resenha, assim como todo trabalho acadêmico, deve ter cunho científico. Ou seja, estar em conformidade com as exigências de objetividade e impessoalidade. Você deve falar apenas do que é mais importante na obra, pois seu leitor raramente estará interessado em muitos detalhes ou em partes menos importantes.

\end{document}